%% kapa: exam=1 round=25 year=2014 period=1 form=A subject=economics subject_ko=경제원론 seq=28 number=68
%% source: Q-Net 감정평가사 기출문제 (cst00309) artlSeq=5101207 file=extracted/1008/5101207/2014년 제25회 감정평가사 1차시험 기출문제/1교시(A형).hwp
%% conversion: kordoc 4.24.51 -> tools/gen1.py (auto); review=proofread-sonnet; flags=none
\begin{question}{68}
총수요 $Y=300+10( \frac{ M } { P } )$과 단기총공급 $Y=500+(P-P^{ e } )$인 경제의 최초균형에서 통화량이 1,000, 물가가 50이다. 통화량이 1,260으로 증가하면 단기균형(ㄱ)과 장기균형(ㄴ)에서 물가는 각각 얼마인가? (단, $Y$는 국민소득, $M$은 통화량, $P$는 물가, $P^{ e }$는 예상물가이다.)
\begin{choices}
  \item ㄱ: 50, ㄴ: 60
  \item ㄱ: 60, ㄴ: 50
  \item ㄱ: 60, ㄴ: 60
  \item ㄱ: 60, ㄴ: 63
  \item ㄱ: 63, ㄴ: 63
\end{choices}
\end{question}
\answer{4}
